% Fig. 3 — vantage points and their blind spots. Used by sections/04-methodology.
% This figure carries the methodological argument: three partial views, one
% ground truth. Uses the shared styles and the \figsub macro from preamble.tex.
\begin{figure}[!t]
\centering
\resizebox{0.94\columnwidth}{!}{%
\begin{tikzpicture}[x=1cm, y=1cm]

  \def\colA{0} \def\colB{3.35} \def\colC{6.7}

  \node[figbox blue, minimum width=2.5cm] (a) at (\colA,0)
    {Introspection\\[1pt]\figsub{names, types}};
  \node[figbox teal, minimum width=2.5cm] (b) at (\colB,0)
    {Traffic capture\\[1pt]\figsub{framing, order}};
  \node[figbox rust, minimum width=2.5cm] (c) at (\colC,0)
    {Application code\\[1pt]\figsub{intent, setup}};

  \node[font=\scriptsize\itshape, text=figslate, align=center] (a2) at (\colA,-1.25)
    {blind to\\preconditions};
  \node[font=\scriptsize\itshape, text=figslate, align=center] (b2) at (\colB,-1.25)
    {blind to\\unused commands};
  \node[font=\scriptsize\itshape, text=figslate, align=center] (c2) at (\colC,-1.25)
    {blind to\\reachability};

  \draw[figarrow] (a) -- (a2);
  \draw[figarrow] (b) -- (b2);
  \draw[figarrow] (c) -- (c2);

  \node[figbox slate, minimum width=8.6cm, minimum height=0.7cm] (t) at (\colB,-2.5)
    {Field telemetry: the only ground truth for execution (P2, H4)};

  \draw[figarrow] (a2) -- (t);
  \draw[figarrow] (b2) -- (t);
  \draw[figarrow] (c2) -- (t);

\end{tikzpicture}%
}
\caption{Each vantage point answers a different question and is blind to the others. A
recovered command is accepted only after state telemetry confirms a physical effect.}
\label{fig:vantage}
\end{figure}
